\chapter{Turn Toward: Menghadap ke Pasangan}

\begin{insightbox}
\textit{``Romansa sejati tidak terjadi di Bali atau Paris. Romansa terjadi di supermarket saat dia bertanya 'Apakah mentega kita habis?' dan kamu menjawab 'Aku cek dulu ya' daripada mengabaikannya.''}
\end{insightbox}

\section{Apa itu ``Turning Toward''?}

\keyterm{Turning Toward} adalah kecenderungan untuk merespons usaha pasangan untuk mendapatkan perhatian, kasih sayang, humor, atau dukungan.

\begin{praktikbox}[Definisi Turning Toward]
\textbf{Turning Toward berarti:}
\begin{itemize}
    \item Menanggapi ``bids'' (ajakan koneksi) dari pasangan
    \item Memberikan perhatian meski sedang sibuk
    \item Menunjukkan bahwa ``kamu penting bagiku''
    \item Membangun \keyterm{Emotional Bank Account} (Tabungan Emosional)
\end{itemize}

\textbf{Contoh sederhana:}
\begin{itemize}
    \item Pasangan: ``Lihat kapal itu besar!''
    \item Turn Toward: ``Wah iya, mirip kapal layar yang kita lihat waktu itu ya?''
    \item Turn Away: (tidak mendengar, terus main HP)
    \item Turn Against: ``Iya, terus? Kenapa?''
\end{itemize}
\end{praktikbox}

\section{Jenis-jenis Bids (Ajakan Koneksi)}

\begin{center}
\renewcommand{\arraystretch}{1.5}
\begin{tabular}{|>{\raggedright}p{3cm}|>{\raggedright}p{5cm}|>{\raggedright\arraybackslash}p{7cm}|}
\hline
\rowcolor{primary!20}
\textbf{Jenis Bid} & \textbf{Contoh} & \textbf{Cara Menanggapi} \\
\hline
\textbf{Attention} & ``Lihat burung itu!'' & Berhenti sejenak, lihat, komentar \\
\hline
\textbf{Affection} & Memegang tangan saat jalan & Sambut, pegang balik, tersenyum \\
\hline
\textbf{Humor} & ``Aku kayak orang gila ya hari ini'' & Tertawa, lempar jokes balik \\
\hline
\textbf{Support} & ``Hari ini berat banget di kantor'' & Dengarkan, validasi, tanya ``Mau cerita?'' \\
\hline
\textbf{Help} & ``Bisa tolong ambilkan air?'' & Bantu dengan sukarela, tanpa komplain \\
\hline
\textbf{Interest} & ``Kamu tadi ngapain?'' & Ceritakan dengan detail, ajak ngobrol \\
\hline
\end{tabular}
\end{center}

\section{Data Gottman yang Mengejutkan}

\begin{warningbox}[Statistik Penting]
Dalam penelitian 6 tahun terhadap pasangan baru menikah:
\begin{itemize}
    \item Pasangan yang \textbf{bertahan}: merespons bids \textbf{86\%} dari waktu
    \item Pasangan yang \textbf{bercerai}: merespons bids hanya \textbf{33\%} dari waktu
    \item \textbf{Selisih hanya 6 respons dari 10 bids} membedakan hubungan yang berhasil vs gagal!
\end{itemize}

\textit{Perbedaan kecil dalam keseharian = Perbedaan besar dalam hasil jangka panjang}
\end{warningbox}

\section{Tiga Respon terhadap Bids}

\begin{center}
\begin{tikzpicture}[
    response/.style={draw=#1, line width=3pt, rounded corners=15pt,
                     fill=#1!10, minimum width=4cm, minimum height=5cm, align=center},
    title/.style={font=\Large\bfseries}
]
    % Turn Toward
    \node[response=sagegreen] (toward) at (0,0) {};
    \node[title, sagegreen] at (0,2) {TURN TOWARD};
    \node[align=left, text width=3.5cm] at (0,0) {
        \begin{itemize}[leftmargin=*, itemsep=0.2cm]
            \item Menanggapi dengan positif
            \item Memberikan perhatian
            \item Menunjukkan minat
            \item Membangun koneksi
        \end{itemize}
    };
    \node[font=\small\bfseries, sagegreen] at (0,-2) {BUILD};
    
    % Turn Away
    \node[response=warmgray] (away) at (6,0) {};
    \node[title, warmgray] at (6,2) {TURN AWAY};
    \node[align=left, text width=3.5cm] at (6,0) {
        \begin{itemize}[leftmargin=*, itemsep=0.2cm]
            \item Mengabaikan
            \item Tidak merespons
            \item Sibuk dengan hal lain
            \item Tidak sadar ada bid
        \end{itemize}
    };
    \node[font=\small\bfseries, warmgray] at (6,-2) {ERODE};
    
    % Turn Against
    \node[response=accent] (against) at (12,0) {};
    \node[title, accent] at (12,2) {TURN AGAINST};
    \node[align=left, text width=3.5cm] at (12,0) {
        \begin{itemize}[leftmargin=*, itemsep=0.2cm]
            \item Merespons dengan negatif
            \item Marah/sarkasme
            \item Menolak kasar
            \item Memulai konflik
        \end{itemize}
    };
    \node[font=\small\bfseries, accent] at (12,-2) {DAMAGE};
\end{tikzpicture}
\end{center}

\section{Emotional Bank Account}

\begin{praktikbox}[Konsep Tabungan Emosional]
Setiap kali kamu \textbf{Turn Toward}, kamu menyetor ke ``Tabungan Emosional''. Setiap kali \textbf{Turn Away}, kamu menarik.

\textbf{Manfaat saldo tinggi:}
\begin{itemize}
    \item Pasangan lebih toleran terhadap kesalahan kecil
    \item Konflik tidak meledak menjadi besar
    \item Rasa percaya lebih kuat
    \item Hubungan tahan terhadap stres eksternal
\end{itemize}

\textbf{Risiko saldo rendah:}
\begin{itemize}
    \item Setiap kesalahan terasa besar
    \item Konflik kecil jadi besar
    \item Negativity Override (cenderung melihat negatif)
    \item Mudah putus asa saat ada masalah
\end{itemize}
\end{praktikbox}

\subsection{50 Activities yang Membangun Tabungan}

\begin{center}
\small
\renewcommand{\arraystretch}{1.2}
\begin{tabular}{|c|l|c|l|}
\hline
\rowcolor{primary!20}
\textbf{No} & \textbf{Aktivitas} & \textbf{No} & \textbf{Aktivitas} \\
\hline
1 & Bertemu di akhir hari, bicara & 26 & Hadiri pertunjukan anak \\
\hline
2 & Belanja bersama & 27 & Pergi ke pesta \\
\hline
3 & Masak/makan malam bersama & 28 & Pergi bersama ke tempat kerja \\
\hline
4 & Bersih-bersih rumah & 29 & Rayakan ulang tahun anak \\
\hline
5 & Belanja kado/pakaian & 30 & Rayakan milestone lain \\
\hline
6 & Makan di luar & 31 & Main game/komputer \\
\hline
7 & Baca/lihat berita bersama & 32 & Atur playdate anak \\
\hline
8 & Bantu self-improvement & 33 & Rencanakan masa depan \\
\hline
9 & Rencanakan dinner party & 34 & Jalan-jalan dengan anjing \\
\hline
10 & Telepon/kirim pesan saat kerja & 35 & Baca bersama \\
\hline
11 & Staycation & 36 & Main board game \\
\hline
12 & Sarapan bersama & 37 & Pentas drama/sandiwara \\
\hline
13 & Pergi ke tempat ibadah & 38 & Jalankan tugas bersama \\
\hline
14 & Kerja kebun/perbaikan rumah & 39 & Lakukan hobi bersama \\
\hline
15 & Kerja relawan & 40 & Ngobrol minum kopi/teh \\
\hline
16 & Olahraga bersama & 41 & Cari waktu bicara tanpa gangguan \\
\hline
17 & Piknik/jalan-jalan & 42 & Gosip (bicara orang) \\
\hline
18 & Kunjungi keluarga & 43 & Hadiri pemakaman \\
\hline
19 & Nonton TV/streaming & 44 & Bantu orang lain \\
\hline
20 & Pesan makanan & 45 & Cari rumah/apartemen \\
\hline
21 & Double date & 46 & Test drive mobil \\
\hline
22 & Baca/bicara di depan api & 47 & Lainnya \\
\hline
23 & Dengarkan musik & 48 & --- \\
\hline
24 & Pergi dancing/konser & 49 & --- \\
\hline
25 & Rayakan ulang tahun anak & 50 & --- \\
\hline
\end{tabular}
\end{center}

\section{Stress-Reducing Conversation}

\begin{tipbox}[Teknik Sakti: The 20-Minute Stress Reducing Conversation]
\textbf{Kapan:} Setiap hari, idealnya di akhir hari\\
\textbf{Berapa lama:} 20-30 menit\\
\textbf{Topik:} Apa saja \textbf{KECUALI} konflik di antara kalian

\textbf{Aturan Emas:}
\begin{enumerate}
    \item \textbf{Don't give unsolicited advice} --- Jangan beri nasihat kecuali diminta!
    \item \textbf{Take turns} --- Bergantian jadi yang cerita
    \item \textbf{Show genuine interest} --- Kontak mata, angguk, ``uh-huh''
    \item \textbf{Communicate understanding} --- Refleksi perasaan
    \item \textbf{Take your partner's side} --- Selalu di pihaknya vs dunia luar
    \item \textbf{Express ``we against others''} --- ``Kita lawan mereka''
    \item \textbf{Show affection} --- Pegang tangan, peluk
    \item \textbf{Validate emotions} --- ``Wajar kamu marah''
\end{enumerate}
\end{tipbox}

\subsection{Contoh Dialog yang Benar}

\begin{praktikbox}[CONTOH: Hari yang Buruk di Kantor]
\textbf{Partner A:} ``Hari ini sial banget. Bosku marah-marah soal deadline, padahal bukan salahku.''

\textbf{Partner B yang BENAR:}
\begin{itemize}
    \item ``Astaga, dia marah-marah? Kok bisa?'' (\textit{genuine interest})
    \item ``Kamu pasti stres banget ya'' (\textit{validation})
    \item ``Memang nggak adil sih kalau begitu'' (\textit{take side})
    \item ``Aku di sini buat kamu'' (\textit{support})
\end{itemize}

\textbf{Partner B yang SALAH:}
\begin{itemize}
    \item ``Yaudah, besok lebih baik'' (\textit{minimize})
    \item ``Mungkin bosmu lagi stres'' (\textit{side with enemy})
    \item ``Kenapa kamu nggak bilang dari awal deadline-nya?'' (\textit{criticism})
\end{itemize}
\end{praktikbox}

\section{Two Obstacles: Mengapa Kita Gagal Turn Toward}

\subsection{Obstacle 1: Bids yang Terselubung}

Kadang bids datang dalam kemasan negatif:

\begin{center}
\renewcommand{\arraystretch}{1.4}
\begin{tabular}{|>{\raggedright}p{6cm}|>{\raggedright\arraybackslash}p{9cm}|}
\hline
\rowcolor{accent!20}
\textbf{Yang Terdengar (Negatif)} & \textbf{Bid Tersembunyi (Positif)} \\
\hline
``Kok nggak pernah beres-beres sih?!'' & ``Aku butuh bantuan. Aku kewalahan.'' \\
\hline
``Kamu cuma mikirin kerjaan!'' & ``Aku kangen. Aku butuh perhatianmu.'' \\
\hline
``Aku capek jadi yang selalu ngajak ngobrol'' & ``Aku butuh usaha darimu juga.'' \\
\hline
``Serah deh, kamu aja yang decide'' & ``Aku butuhmu peduli sama pendapatku.'' \\
\hline
\end{tabular}
\end{center}

\begin{tipbox}[Cara Melihat Bid Tersembunyi]
\begin{enumerate}
    \item \textbf{Pause:} Jangan reaksi langsung saat diserang
    \item \textbf{Ask:} ``Kamu butuh aku untuk apa sekarang?''
    \item \textbf{Look beneath:} Apa kebutuhan emosional di balik kata-kata kasar?
    \item \textbf{Respond to bid:} Jawab kebutuhan, bukan kata-kata
\end{enumerate}
\end{tipbox}

\subsection{Obstacle 2: Digital Distraction}

\begin{warningbox}[The Wired World Problem]
\begin{itemize}
    \item Rata-rata orang cek HP 96 kali per hari
    \item Notifikasi terus-menerus memecah perhatian
    \item ``Phubbing'' (phone + snubbing) membunuh koneksi
    \item Digital distraction = Turning Away tersembunyi
\end{itemize}
\end{warningbox}

\textbf{Solusi:}
\begin{enumerate}
    \item Tetapkan ``sacred time'' tanpa gadget (misal: makan malam)
    \item Buat ``phone hotel'' --- tempat meletakkan HP saat di rumah
    \item Notifikasi off untuk aplikasi non-esensial saat bersama
    \item Ritual: HP di-charge di luar kamar tidur
\end{enumerate}

\section{Exercise: Emotional Bank Account Ledger}

\begin{exercisebox}[Cara Menggunakan Ledger]
\textbf{Setiap hari, catat:}
\begin{enumerate}
    \item Setiap kali pasangan ``menyetor'' (Turn Toward)
    \item Setiap kali pasangan ``menarik'' (Turn Away)
    \item Setiap kali kamu menyetor/menarik
\end{enumerate}

\textbf{Akhir minggu:}
\begin{itemize}
    \item Review pola
    \item Diskusikan dengan pasangan
    \item Komitmen untuk minggu depan
\end{itemize}

\textbf{Ingat:} Bukan untuk kompetisi, tapi untuk kesadaran!
\end{exercisebox}

\section{When Your Spouse Doesn't Turn Toward}

Jika kamu merasa sering ditolak/diabaikan:

\begin{enumerate}
    \item \textbf{Jangan personalisasi} --- Mungkin dia sedang stres/bukan tentangmu
    \item \textbf{Check timing} --- Apakah timing bid-mu tepat?
    \item \textbf{Soft start-up} --- Coba cara berbeda untuk meminta koneksi
    \item \textbf{Talk it out} --- Diskusikan pola ini saat waktu tenang
    \item \textbf{Give benefit of the doubt} --- Asumsikan niat baik
\end{enumerate}

\begin{exercisebox}[Exercise: Talking It Out]
Jika salah satu merasa sering ditolak, isi form ini dan diskusikan:

\textbf{Saat ini terjadi, aku merasa:}
\begin{itemize}
    \item[$\square$] Dikesampingkan
    \item[$\square$] Tidak dihargai
    \item[$\square$] Kesepian
    \item[$\square$] Marah
    \item[$\square$] Salah paham
\end{itemize}

\textbf{Yang memicu perasaan ini:}
\begin{itemize}
    \item[$\square$] Aku merasa dikesampingkan
    \item[$\square$] Aku butuh lebih banyak perhatian
    \item[$\square$] Aku merasa tidak diinginkan
\end{itemize}

\textbf{Ini mengingatkanku pada:}
\begin{itemize}
    \item[$\square$] Cara aku diperlakukan di keluarga
    \item[$\square$] Hubungan sebelumnya
    \item[$\square$] Trauma/kecederaan lama
    \item[$\square$] Harapan yang tidak terpenuhi
\end{itemize}

\textbf{Kontribusiku pada masalah ini:}
\begin{itemize}
    \item[$\square$] Aku terlalu sensitif
    \item[$\square$] Aku tidak mengekspresikan dengan jelas
    \item[$\square$] Aku sering sibuk juga
\end{itemize}
\end{exercisebox}
